\subsection{A uniform upper bound for the original cusp norms}
\label{shared-bergman-actual-upper}

We retain the integral model, its complex uniformization and its
Hermitian metric.  The probability measure on a component is the one
induced by the original adelic measure~\cite{bergman-upper-source}.
In particular, normalizing a
component does not permit us to discard its arithmetic volume.

Let $E$ be the original CM field, let $F=E^+$, and put
\[
 f=[F:\mathbb Q],\qquad \Delta=|\operatorname{disc}E|,
 \qquad \operatorname{rd}(E)=\Delta^{1/(2f)}.
\]
The original hypotheses include that $E/F$ is unramified at all finite
primes.  The rank-$r$, $r\ge2$, Hermitian lattice used in the principal-level
description is self-dual over $\mathcal O_E$; it need not be free.
Write $M\ge3$ for the level and $d_H\ge1$ for the original positive
quadratic parameter.  At an orthogonal frame $x$, write
\[
 q_x(v)=\sum_{\tau\in\Psi}\operatorname{maj}_{x,\tau}(v)
\]
for the sum of the positive Hermitian majorants, where $\Psi$ contains
one embedding from each conjugate pair.  If $\lambda_{\tau,i}$ are
the model weights, assume
\[
 \sum_{\tau\in\Psi}
 \bigl(\max_i\lambda_{\tau,i}-\min_i\lambda_{\tau,i}\bigr)
 \le w,\qquad w\ge1.
\]
The weights in this statement are the actual model weights after the
prescribed signed permutation.  Let $\mathcal H_Z$ be the original
finite-dimensional cusp section space on a component $Z$, with its
component probability norm $\|\cdot\|_{2,Z}$.  Its Bergman function is
\[
 B_Z(x)=\sup_{0\ne s\in\mathcal H_Z}
              \frac{\|s(x)\|^2}{\|s\|_{2,Z}^2},
\]
with value zero when $\mathcal H_Z=0$.

Set
\begin{align}
 D&=fr(r-1)/2,&
 \rho&=\frac1{200r^3(w+f)},\label{eq:actual-upper-radius}\\
 \epsilon&=\min\left\{1,
 \frac{\pi}{2M\sqrt{d_H}(fw+2)\log2}\right\},&
 A_\epsilon&=
 \left(1+\frac{2(2fr)^2\Delta}{\epsilon}\right)^{2fr^2},
 \label{eq:actual-upper-multiplicity}\\
 V&=2h_E M^{2f(r^2+1)}
       \operatorname{rd}(E)^{f(r^2-1)/2},&
 C_{\mathrm{sup}}&=
       \frac{8D!}{(\pi\rho^2)^D}A_\epsilon V,
 \label{eq:actual-upper-constant}
\end{align}
where $h_E$ denotes the ideal class number of $E$.

\begin{theorem}[The original upper Bergman bound]
\label{thm:actual-upper-bergman}
For the original good principal-level cover, every coefficient
embedding, every component $Z$ and every $s\in\mathcal H_Z$ satisfy
\[
 \sup_{x\in Z}\|s(x)\|^2
       \le C_{\mathrm{sup}}\|s\|_{2,Z}^2,
 \qquad B_Z(x)\le C_{\mathrm{sup}},
 \qquad \dim\mathcal H_Z\le C_{\mathrm{sup}}.
\]
The statement uses the classical Tamagawa-number theorem for special
unitary groups.  All comparisons from its invariant differential
measure to the displayed constant are given below.  No bound is
asserted here for a ramified extension $E/F$ or a different integral
lattice.

If $L\ge1$ and
\begin{gather}
 \log(fr)\le L,\quad \log M\le30L,
 \quad \log\Delta\le4fL,\quad
 \log d_H\le4L,\quad \log w\le2L,
 \label{eq:actual-upper-size}
\end{gather}
then
\begin{equation}\label{eq:actual-upper-logbound}
            1\le C_{\mathrm{sup}},\qquad
            \log C_{\mathrm{sup}}\le250f^2r^2L.
\end{equation}
In the original family, let $L_b\ge1$.  Suppose
\[
 \log M\le30L_b,\qquad
 \log(2f)\le3L_b,\qquad \log\operatorname{rd}(E)\le3L_b,\qquad
 \log(fr)\le L_b,
\]
and put $S=\sum_{\tau,i}|u_{\tau,i}|$.  If $S=0$ or
$\log S\le20L_b$, then one may take
\begin{equation}\label{eq:actual-upper-family}
                   B_Z(x)\le \exp\bigl(4000(frL_b)^3\bigr).
\end{equation}
Here $u$ is the original integral weight and the original host
discriminant relation $d_H^f\le\Delta$ is retained.
\end{theorem}

\begin{proof}
\emph{A half-maximum point has a lower lattice minimum.}
The zero section requires no argument.  For a nonzero section choose
a frame $x$ with point norm greater than half its supremum, then
orthogonalize the frame.  The original norm is unchanged.  The
integral Hermitian lattice $L'$ and principal-level transvections
constructed in the short-theta proof give
\begin{equation}\label{eq:actual-upper-minimum}
                    q_x(v)\ge\epsilon
                  \qquad(0\ne v\in L').
\end{equation}
We recall why this estimate has precisely the displayed level
dependence.  If all embedding majorants of $v$ are less than
$\epsilon\le1$, its integral Hermitian norm has all conjugates of
absolute value less than one and is therefore zero.  Its rank-one
unipotent orbit, with complex parameter $\zeta$, remains in the
domain for $\Im\zeta>-T$, where
$T=(k_0\epsilon)^{-1}$.  The period $M_*$ of the original
principal-level action obeys
\[
                       M_*k_0=2M\sqrt{d_H}.
\]
The actual arithmetic character is trivial on this action, so the
evaluated section is periodic with period $M_*$.  Its constant
Fourier coefficient is zero by cuspidality.  On the line
$\Im\zeta=-T/2$, its norm differs from the norm at $x$ by at most
$2^{w/2}$ in the frame factor.  Applying the disk Schwarz lemma
after $q=\exp(2\pi i\zeta/M_*)$ gives
\[
              1\le2\,2^{w/2}\exp(-\pi T/M_*).
\]
But the definition of $\epsilon$ makes the right side at most
$2^{1+w/2-(fw+2)}<1$.  This contradiction proves
\eqref{eq:actual-upper-minimum}.  The arithmetic character assertion
is the one for the constructed bundle: its character law on all
nonnegative antitone exponent arrays extends, by differences, to
all integral model weights.  Finite image alone would not suffice.

\emph{The actual chart has bounded multiplicity.}
Put $n=2fr$.  The real quadratic form $q_x$ has dimension $n$.
For its Euclidean dual lattice $(L')^*$, the trace-pairing identity
and Hermitian self-duality imply
\[
             \frac{\Delta}{2}J_x(L')^*\subseteq L',
\]
where $J_x$ is the orthogonal involution of the Hermitian majorant.
Indeed, trace duality gives the inverse different, the integer
$\Delta$ clears that different, and the sum over a CM type is
half the full trace.  Thus the dual minimum is at least
$2\sqrt\epsilon/\Delta$.

The elementary transference argument already proved for the short
theta span gives
$\lambda_n(L')\lambda_1((L')^*)\le n^2$.
Consequently $L'$ contains $r$ vectors $v_1,\ldots,v_r$ independent
over $E$, with
\begin{equation}\label{eq:actual-upper-short-basis}
                  q_x(v_i)\le\frac{n^4\Delta^2}{4\epsilon}.
\end{equation}
The argument first gives $n$ vectors independent over $\mathbb Q$;
one then selects an $E$-basis.  It does not select a free
$\mathcal O_E$-basis or change the integral lattice.

Use the original lower-triangular flag chart of radius $\rho$ at
$x$.  Comparing the two positive/negative orthogonal projections
in its Gram matrices gives
\[
 \tfrac12q_x(v)\le q_z(v)\le2q_x(v)
       \qquad(\|z\|\le\rho).
\]
For clarity, write the chart matrix as $I+N$ in an orthogonal frame.
Its entries satisfy $|N_{ab}|\le\rho$ and
$\|N\|_{\mathrm{op}}\le r\rho$; the Gram perturbation has norm
at most $2r\rho+r^2\rho^2$.  The positive block and its inverse
remain within the same bounds by the Neumann series.  In the
formula for the positive projection this gives an operator norm
error at most $8r^2\rho<1/2$.  This proves the displayed
comparison, for the indefinite signatures as well as for the
definite endpoints.

If an arithmetic element $\gamma$ returns this chart to itself,
the intermediate upper-triangular change of flag leaves the
majorant invariant.  Hence
$q_x(\gamma v)\le4q_x(v)$.  The map
\[
               \gamma\longmapsto(\gamma v_1,\ldots,\gamma v_r)
\]
is injective.  Balls of radius $\sqrt\epsilon/2$ around lattice
points are disjoint.  Comparing their volumes shows that a ball
of radius $R$ contains at most
$(1+2R/\sqrt\epsilon)^n$ lattice points.  Taking
$R=n^2\Delta/\sqrt\epsilon$ in
\eqref{eq:actual-upper-short-basis}, and then taking the $r$-th
power, bounds the number of returning elements by $A_\epsilon$.

\emph{A local mean inequality retains the metric and measure.}
In the same chart let $\delta_{\tau,j}$ be the leading Hermitian
Gram determinants, with $\delta_{\tau,0}=1$.
The triangular chart has determinant one, so
$|\delta_{\tau,r}|=1$.  The same matrix estimates give
\[
        \bigl|\log|\delta_{\tau,j}|\bigr|\le9r^2\rho.
\]
In the exact metric formula, subtracting the midpoint of the
largest and smallest weights costs nothing because the two end
determinants have absolute value one.  Summation by parts therefore
gives a squared frame factor at least $\exp(-1/20)$.
The invariant flag density is the product of
$|\delta_{\tau,j}|^{-2}$, and is at least $\exp(-1/10)$.
Both assertions follow directly from
$200r^3\rho(w+f)=1$.  Their product exceeds $1/2$.

The mean-value inequality for the squared absolute value of the
holomorphic coefficient, on a ball in $\mathbb C^D$, now gives
\[
 \|s(x)\|^2\le
 \frac{2D!}{(\pi\rho^2)^D}
           \int_{\|z\|<\rho}\|s(z)\|^2\,d\mu_{\mathrm{flag}}(z).
\]
Unfolding this nonnegative integral to a fundamental set multiplies
it by at most $A_\epsilon$.  The half-maximum choice costs a factor
four.  Thus, on this uniformizing component,
\begin{equation}\label{eq:actual-upper-raw}
 \sup\|s\|^2\le
 \frac{8D!}{(\pi\rho^2)^D}A_\epsilon
       \int_{\mathcal F}\|s\|^2\,d\mu_{\mathrm{flag}}.
\end{equation}
Tonelli's theorem applies throughout.  The argument needs neither
an injective chart nor an unspecified compactness constant.

\emph{The arithmetic volume in the probability normalization.}
The required bound for the sum of the original fundamental-set
volumes is
\begin{equation}\label{eq:actual-upper-mass}
             \sum_k\mu_{\mathrm{flag}}(\mathcal F_k)\le V.
\end{equation}
Here is the normalization calculation.  At a finite place of $F$
with residue size $q$, the self-dual lattice in the unramified
quadratic algebra defines a smooth special unitary group scheme.
Its volume for an integral invariant differential is
\[
 a_v=\prod_{j=2}^r(1-\chi_v^j q^{-j}),
 \qquad
 \chi_v=\begin{cases}1&\text{at a split place},\\-1&\text{at an inert place}.
                    \end{cases}
\]
The formula follows from the orders of the split and unitary
finite groups and smooth lifting.  It holds at primes over two
and at primes dividing $r$; no division by these primes is used.
It also gives
\[
                 \prod_v a_v^{-1}
                   \le\prod_{j=2}^r\zeta_F(j)
                   \le\left(\prod_{j=2}^r\zeta(j)\right)^f.
\]

The only global volume theorem used is that the Tamagawa number
of the simply connected group $\operatorname{SU}(h)$ is one.
For example, this is the outer type-$A$ case recorded in
Prasad~\cite[\S3.3]{bergman-upper-prasad}; his \S3.7 gives the corresponding local-factor
formula.  Summing over all finite double cosets gives the same
identity when all real factors are compact, so no isotropy or
strong-approximation assumption is imposed at this step.

To compare its differential measure with the original Riemannian
measure, choose any integral $F$-basis of the trace-zero
anti-Hermitian Lie algebra, spanning a sublattice of index $a$.
The trace form on the full endomorphism lattice is unimodular
locally because $E/F$ is unramified and the Hermitian lattice is
self-dual.  Passing to trace zero contributes the determinant
$r$: in diagonal coordinates its Gram matrix is
$I_{r-1}+\mathbf1\mathbf1^t$, of determinant $r$.
The off-diagonal pairs have unit determinant.  This computation
may be made after unramified faithfully flat extension and hence
gives the same determinant ideal before extension.  It proves
that the real volume conversion is $r^{f/2}a$.
The same $a$ occurs inversely in the finite differential volumes,
so it cancels.  The sum of the Riemannian covolumes at level one is
therefore
\[
        r^{f/2}|\operatorname{disc}F|^{(r^2-1)/2}
                             \prod_v a_v^{-1}.
\]
This computation allows a projective integral Lie lattice and
does not assume it has a global integral basis.

The compact diagonal torus has Riemannian volume
$(2\pi)^{r-1}\sqrt r$ at each real place; the determinant-one
condition gives the Gram matrix
$I_{r-1}+\mathbf1\mathbf1^t$.  Each complex off-diagonal flag
coordinate contributes the factor two.  Consequently the exact
conversion from a flag fundamental-set volume to its special
unitary covolume, for a torsion-free arithmetic group, is
\[
       \left(2^{r(r-1)/2}(2\pi)^{r-1}\sqrt r\right)^f.
\]
Since $\zeta(j)<2$ for $j\ge2$, the ratio of the sum of
level-one group covolumes to this factor is at most
$|\operatorname{disc}F|^{(r^2-1)/2}$.  At level one this is a
volume with stabilizer weights.  We use the ordinary flag
fundamental-set volume only after passing to the principal level.

It remains to retain the determinant classes and level.
An element of the principal congruence group has determinant
$a\in\mathcal O_E^\times$ with $a\bar a=1$ and
$a-1\in M\mathcal O_E$.  If $a\ne1$, the norm of $a-1$ is
both at most $2^{2f}$ and at least $M^{2f}$, a contradiction.
Thus these rational stabilizers really lie in the special unitary
group; this is where $M\ge3$ is used.
The principal group is also torsion-free.  Represent it on the
underlying integral lattice.  If an odd prime $p$ divides $M$,
the group $1+p\operatorname{Mat}_n(\mathbb Z_p)$ is torsion-free;
if $M$ has no odd prime factor then $4\mid M$, and one uses
$1+4\operatorname{Mat}_n(\mathbb Z_2)$.
For completeness, if $A-1=p^aB$ has least entry valuation $a$,
the binomial expansion shows that $A^p-1$ has least entry
valuation $a+1$ in these ranges.  Reduction shows that any
finite order would be a power of $p$, which is impossible.
The action on the flag domain is therefore free, since a
discrete subgroup meets its compact point stabilizer in a
finite group.  This justifies the preceding measure conversion
at the actual level, including the central stabilizers.
Hilbert 90 writes a finite norm-one idele as $b/\bar b$.
Choose the ideal class of $b$ among $h_E$ representatives; the
remaining integral units, modulo $M$, have at most $M^{2f}$
residue classes.  Their quotient has norm one and is congruent
to one modulo $M$.  Locally it is the determinant of an integral
unitary transformation congruent to one: diagonalize the
unramified unimodular Hermitian lattice and act on one line.
Descent gives the assertion for its original lattice.
Thus the number of determinant classes is at most
$h_E M^{2f}$.  In each determinant class the principal level
index is at most $M^{2fr^2}$, by its injection into endomorphisms
of the rank-$r$ module modulo $M$.  These two bounds, with the
harmless factor two retained in $V$, prove
\eqref{eq:actual-upper-mass}.  Finally,
$\Delta=|\operatorname{disc}F|^2$ follows from finite
unramifiedness, giving the exponent in
\eqref{eq:actual-upper-constant}.

The constructed component dictionary identifies its probability
norm with the sum of these integrals divided by the mass of the
corresponding components.  That mass is at most the full sum in
\eqref{eq:actual-upper-mass}.  Applying
\eqref{eq:actual-upper-raw} and then this exact normalization proves
the asserted pointwise norm inequality.  Taking its supremum over
unit sections proves the Bergman bound.  Integrating an orthonormal
basis gives $\int_Z B_Z\,d\mu_Z=\dim\mathcal H_Z$, and proves the
dimension bound.

\emph{The numerical estimate.}
Minkowski's ideal-class theorem gives an integral ideal of norm
at most $\sqrt\Delta$ in each class.  Indeed its CM constant is
at most $(2/\pi)^f<1$, by bounding the first $f$ factors in
$(2f)!/(2f)^{2f}$ by $1/2$.  Counting ideals with the Dirichlet
series at two gives
\[
                 h_E\le\Delta\zeta_E(2)
                        \le4^f\Delta.
\]
Under \eqref{eq:actual-upper-size} this implies
$\log V\le80fr^2L$.
Also
\[
       \epsilon^{-1}\le3fwM\sqrt{d_H},\qquad
       \log\epsilon^{-1}\le2+35L.
\]
Using $1+2n^2\Delta/\epsilon\le3n^2\Delta/\epsilon$
and $n=2fr$, we obtain
\[
          \log A_\epsilon
                  \le86fr^2L+8f^2r^2L.
\]
Since $D!\le D^D$, $D\le fr^2/2$, $\log D\le2L$,
$\log(w+f)\le\log2+2L$ and $\log2<3/4$, the chart factor obeys
\[
        \log\frac{8D!}{(\pi\rho^2)^D}\le14fr^2L.
\]
Adding gives at most $188f^2r^2L$, and hence the stated bound
$250f^2r^2L$.  Each factor in $C_{\mathrm{sup}}$ is at least
one, so its lower bound also holds.

For the family statement put $S=\sum_{\tau,i}|u_{\tau,i}|$ and
$w=4fS+1$.  The exact model-weight formula gives total spread at
most $4fS$.  If $S=0$, then $w=1$.  Otherwise $S\ge1$ because
the weights are integral; the displayed hypotheses imply
$\log w\le25L_b$.  The root-discriminant bound gives
$\log\Delta\le6fL_b$, and $d_H^f\le\Delta$ gives
$\log d_H\le6L_b$.  Apply
\eqref{eq:actual-upper-logbound} with $L=25L_b/2$.
It gives
\[
      \log C_{\mathrm{sup}}\le3125f^2r^2L_b
                                 \le4000(frL_b)^3.
\]
All constants are independent of the component, embedding and
growing arithmetic field.
\end{proof}

\begin{corollary}[The multiplication norm for the original sections]
\label{cor:actual-upper-multiplication}
If $s$ has the first weight and $t$ the second weight in the
original section multiplication, then
\[
       \|st\|_{2,Z}\le\sqrt{C_{\mathrm{sup}}}
                              \|s\|_{2,Z}\|t\|_{2,Z}.
\]
This holds for the original integral multiplication map whenever
$s,t$ are integral sections, since their pointwise metrics multiply.
\end{corollary}
\begin{proof}
Integrate $\|s(x)t(x)\|^2\le
(\sup\|s\|^2)\|t(x)\|^2$ against the same component probability
measure and apply Theorem~\ref{thm:actual-upper-bergman}.
\end{proof}

The theorem supplies the global upper Bergman, supremum and rank
estimates with their original numerical family bound, and the
corollary supplies the corresponding multiplication norm.  The
thin-region lower bound, point-height comparison and final arithmetic
section remainder remain separate.  Generation of the entire
integral cusp lattice from the short coefficient-field spanning
family is treated in Section~\ref{shared-actual-integral-cusp-frame}.
